\maketitle
\head{本版编排与收录说明}
\noindent\textbf{这是一份按内容顺序整理的中文工作译稿，不是全部上传材料的完结全译本。}
第I、II章及第III章的两个已出版分册正文已经连续收入；第0章当前连续收入
§§1--23，止于（23.2.10，ii）及原著“待续”。
第IV章主体及第V卷预备笔记余文
尚未连续译入。未译部分没有用原文扫描页、摘要或空白标题替代译文。

\section*{正文次序}
正文依原著内容编号排列：导言；第0章预备知识；第I章概形的语言；
第II章若干类态射的初等整体研究；第III章凝聚层的上同调研究；
后续待译范围说明；独立编排的第V卷预备笔记已译部分；
最后为各原分册书末资料。
第III章第一部分的§§1--5与第二部分的§§6--7直接衔接。

第0章原先分载于第一卷、第三卷第一部分、第四卷第一部分。
合订时将已译补篇归入同一章，仍保留
$0_{\mathrm I}$、$0_{\mathrm{III}}$、$0_{\mathrm{IV}}$ 等出处记号。
这些是原分册标识，不是三个重复的章；按编号排序也不表示所有跨章前提
都已经在前文出现。原著交叉引用仍保留原号。

\section*{实际收录范围}
\begin{longtable}{@{}p{.25\linewidth}p{.67\linewidth}@{}}
\toprule 材料&已收入的连续译文及状态\\\midrule\endhead
第一卷（1960）&导言；第0章§§1--7；第I章§§1--10，止于（10.15.7）；相关书末资料。沿用既有校订源码。\\
第二卷（1961）&第II章§§1--8，止于（8.14.14）；参考文献、索引、目录及勘误清单一。书末资料沿用此前依底本恢复的译排来源。\\
第三卷第一部分（1961）&第0章§§8--13；第III章§§1--5，止于（5.5.2）；相关书末资料。\\
第三卷第二部分（1963）&第III章§§6--7，止于（7.9.14），及证明、记号索引、术语索引、原目录、勘误清单二。\\
第四卷第一部分（1964）&第0章补篇章首、§§14--23，止于（23.2.10，ii）末尾。旧条目（19.7.1.5）的较弱原版拓扑范围仍待核，其原文与已核定充分形式的差别保存在脚注及报告中。\\
第V卷英文预备笔记&\texttt{EGAV1-2.ps}六页可读材料；\texttt{EGAVaiello.pdf}卷首说明及“超平面截面与锥投影”§§1--6。余文不计入已译范围。\\
\bottomrule
\end{longtable}

\section*{正文、脚注与核查报告}
新增正文采用已经核实的正确陈述、公式与证明；原文的错误写法、修改理由
及依据置于脚注。作者正式勘误与译者补证分开，较长反例及计算收入独立报告。
未确认的疑点不擅自定为错误；原著本来省略的证明及预备笔记中缺文、不可辨认
之处，不凭通常知识补造成原文。本轮新注号以P16开头，既有注号保留。

这一方式适用于本轮续译范围，并承接第七至十五版列明的回溯修订。
\textbf{它不表示全部旧稿中的历史核查意见已经逐项改入正文。}
本版保存旧译文，不据此次编译声称其每一句或每个证明均经重新审定。
作者撤销的（19.3.12）在正文仍只保留编号与撤销状态；
完整历史陈述与证明保存于报告。

第V卷上传载体是预备笔记的翻译与编辑，英译编者为Piotr Blass与Joseph Blass。
其中的重组建议、未完成句子及待补证明，不与正式出版的第I--IV章标为同一状态。
与合并英文稿重合的PS材料不重复计算。

\section*{第十六版新增范围}
本轮连续译入第0章§23“日本环”全文：23.1日本环与23.2诺特局部整环的整闭包。
范围从（23.1.1）至（23.2.10，ii）及原著“待续”，包括两个子编号引理，共19个原著编号条目、3个编号公式。
底本为第四卷第一部分原分册印刷第213--221页的相关部分（上传PDF第212--220页）。
至此现有底本中第0章§§1--23的正文连续收入；下一接续点是第IV章内容提要、章首说明及§1（1.1.1）。

（23.2.7）的有限中间环改为半局部环；（23.2.7.1）补出首一方程越过幂零根的提升及忠实平坦下降。
Tate定理应用的完备性、素理想高度与清分母步骤作必要补明。
作者关于（23.1.3）的两项正式勘误落实在正文；其余编号、引文和识读问题分开登记。
本轮未给出Tate、Nagata及Mori主要结论的反例，不将辅助断言和证明步骤的问题混称为核心定理错误。

此前数学正文和历史报告保留，不据本轮编译声称重新审定旧译文。
旧版H8-P01等事项继续保留；第IV章主体和第V卷预备笔记余文尚未译完。

\section*{版式与使用}
使用 \texttt{\string\documentclass[UTF8,a4paper]\{ctexbook\}}，
由文档类决定默认字体、字号、版心、标题、页眉与双面分页。
不另设页边距，不自定义页眉；保留原章号、条目号和公式号。目录正文页码进入四位数后，仅将目录页码预留宽度增至2.4em（相应右侧预留3.4em），避免数字溢出；正文字体与版心不变。
默认双面排版所需的空白偶数页不删除。
前置部分用罗马数字页码，正文及书末资料连续编号；
总目录指向合订本页面，原著目录所列页码仍指原版。

本版同时提供PDF、可独立编译的单文件TeX和分章节工程。
实际编译、结构与渲染检查结果见工程记录。
这些检查不代替出版级逐句审校，也不是全部数学证明的独立形式化验证。
\tableofcontents
